\problem{4.21}{Ladder normalization for orbital angular momentum}
With $L_\pm=L_x\pm\ii L_y$ and
$L_\pm^\dagger=L_\mp$,
\begin{equation}
L_\mp L_\pm=L^2-L_z^2\mp\sh L_z.
\end{equation}
If $\ket{\ell m}$ is normalized,
$\|L_+\ket{\ell m}\|^2=
\bra{\ell m}L_-L_+\ket{\ell m}
=\sh^2[\ell(\ell+1)-m(m+1)]$.
Therefore
\begin{align}
\boxed{L_+\ket{\ell m}}
&=\boxed{\sh\sqrt{(\ell-m)(\ell+m+1)}
\ket{\ell,m+1}},\\
\boxed{L_-\ket{\ell m}}
&=\boxed{\sh\sqrt{(\ell+m)(\ell-m+1)}
\ket{\ell,m-1}}.
\end{align}
The top and bottom coefficients vanish at $m=\pm\ell$,
respectively.

\problem{4.22}{Orbital commutators and rotational invariance}
Since $L_z=xp_y-yp_x$,
\begin{align}
[L_z,x]&=[xp_y,x]-[yp_x,x]=\ii\sh y,&
[L_z,y]&=-\ii\sh x,&[L_z,z]&=0,\\
[L_z,p_x]&=\ii\sh p_y,&
[L_z,p_y]&=-\ii\sh p_x,&[L_z,p_z]&=0.
\end{align}
Insert $L_x=yp_z-zp_y$ into
$[L_z,L_x]$ and use the product rule
to obtain $\boxed{[L_z,L_x]=\ii\sh L_y}$.
The same commutators imply
$[L_z,x^2+y^2+z^2]=0$ and
$[L_z,p_x^2+p_y^2+p_z^2]=0$;
hence
$\boxed{[L_z,r]=[L_z,p^2]=0}$.
For central $V(r)$, $[L_i,V(r)]=0$
and $[L_i,p^2]=0$ for all $i$,
so $\boxed{[H,\mathbf L]=0}$ and
$[H,L^2]=0$.

\problem{4.23}{Quantum torque and orbital-angular-momentum conservation}
For $L_i=\epsilon_{ijk}r_jp_k$,
apply the generalized Ehrenfest theorem:
\begin{align}
\frac{\dd}{\dd t}\avg{L_i}
&=\frac{\ii}{\sh}
\epsilon_{ijk}\avg{[H,r_jp_k]}\\
&=\epsilon_{ijk}\avg{(p_j/m)p_k
-r_j\partial_kV}\\
&=-\epsilon_{ijk}\avg{r_j\partial_kV}.
\end{align}
The term with $p_jp_k$ vanishes
because it is symmetric in $j,k$.
Thus
\begin{equation}
\boxed{\frac{\dd\avg{\mathbf L}}{\dd t}
=\avg{\mathbf r\times(-\nabla V)}}.
\end{equation}
For $V=V(r)$, $\nabla V=V'(r)\hat{\mathbf r}$,
so the torque vanishes identically and
$\boxed{\dd\avg{\mathbf L}/\dd t=0}$.

\problem{4.24}{Differential representation of $L^2$}
For a test function $f(\theta,\phi)$,
\begin{equation}
L_\pm f=\pm\sh\ee^{\pm\ii\phi}
(\partial_\theta\pm\ii\cot\theta\,\partial_\phi)f.
\end{equation}
Apply $L_+$ to $L_-f$, differentiating
both the exponential and $\cot\theta$:
\begin{equation}
L_+L_-=-\sh^2\left(
\partial_\theta^2+\cot\theta\,\partial_\theta+
\cot^2\theta\,\partial_\phi^2+
\ii\partial_\phi\right).
\end{equation}
Using the algebraic identity
$L^2=L_+L_-+L_z^2-\sh L_z$
and $L_z=-\ii\sh\partial_\phi$,
the first-derivative terms cancel.
The coefficient of $\partial_\phi^2$
combines as $1+\cot^2\theta=\csc^2\theta$:
\begin{equation}
\boxed{L^2=-\sh^2
\left[\frac1{\sin\theta}\partial_\theta
(\sin\theta\,\partial_\theta)+
\frac1{\sin^2\theta}\partial_\phi^2\right]}.
\end{equation}

\problem{4.25}{Highest-weight spherical harmonic}
By the ladder algebra,
$\boxed{L_+Y_\ell^\ell=0}$.
Let $Y_\ell^\ell(\theta,\phi)=
F(\theta)\ee^{\ii\ell\phi}$.
Then $L_+Y=0$ reduces to
$F'-\ell\cot\theta\,F=0$ and hence
$F(\theta)=C\sin^\ell\theta$.
Normalize:
\begin{equation}
1=2\pi|C|^2\int_0^\pi
\sin^{2\ell+1}\theta\dd\theta,\quad
\int_0^\pi\sin^{2\ell+1}\theta\dd\theta
=\frac{2^{2\ell+1}(\ell!)^2}{(2\ell+1)!}.
\end{equation}
Choosing the standard phase gives
\begin{equation}
\boxed{Y_\ell^\ell=
(-1)^\ell\frac{\sqrt{(2\ell+1)!}}
{2^\ell\ell!\sqrt{4\pi}}\,
\sin^\ell\theta\ee^{\ii\ell\phi}}.
\end{equation}
This agrees with the associated-Legendre formula.

\problem{4.26}{Raising $Y_2^1$ to obtain $Y_2^2$}
Start from
$Y_2^1=-\sqrt{15/(8\pi)}
\sin\theta\cos\theta\,\ee^{\ii\phi}$.
Apply
$L_+=\sh\ee^{\ii\phi}
(\partial_\theta+\ii\cot\theta\,\partial_\phi)$:
\begin{equation}
L_+Y_2^1=
\sh\sqrt{\frac{15}{8\pi}}
\sin^2\theta\,\ee^{2\ii\phi}.
\end{equation}
The normalization factor from Problem 4.21
is $\sh\sqrt{(2-1)(2+1+1)}=2\sh$.
Therefore
\begin{equation}
\boxed{Y_2^2=
\sqrt{\frac{15}{32\pi}}\,
\sin^2\theta\,\ee^{2\ii\phi}}.
\end{equation}

\problem{4.27}{Rigid rotor spectrum of a diatomic molecule}
The center-of-mass distances are
$r_1=m_2a/(m_1+m_2)$ and
$r_2=m_1a/(m_1+m_2)$, so
\begin{equation}
I=m_1r_1^2+m_2r_2^2
=\boxed{\frac{m_1m_2}{m_1+m_2}a^2}.
\end{equation}
For a rigid rotor $H=L^2/(2I)$, giving
\begin{equation}
\boxed{E_\ell=\frac{\sh^2\ell(\ell+1)}{2I}},
\quad\boxed{\psi_{\ell m}=Y_\ell^m(\theta,\phi)},
\quad d_\ell=2\ell+1.
\end{equation}
For electric-dipole rotational transitions
$\Delta\ell=\pm1$, the emission frequencies are
\begin{equation}
\boxed{\nu_j=\frac{E_j-E_{j-1}}h
=\frac{\sh j}{2\pi I}},\quad j=1,2,\dots.
\end{equation}
Adjacent lines are separated by
$\Delta\nu=\sh/(2\pi I)$.
Using the approximate separation in the
textbook CO figure, $\Delta\tilde\nu\simeq
3.84\ {\rm cm^{-1}}$,
$I=\sh/(2\pi c\Delta\tilde\nu)$.
The $^{12}$C--$^{16}$O reduced mass is
about $6.86$ atomic mass units, giving
\begin{equation}
\boxed{a=\sqrt{I/\mu}\simeq1.13\times10^{-10}\ {\rm m}}.
\end{equation}
\textit{The line spacing is read approximately from a figure;
the quoted bond length is correspondingly approximate.}

\problem{4.28}{Why a classical spinning electron cannot explain spin}
For a uniform classical solid sphere,
$I=(2/5)m_er_c^2$, where
$r_c=e^2/(4\pi\epsilon_0m_ec^2)$.
Equating $I\omega=\sh/2$ gives
\begin{equation}
v_{\rm equator}=\omega r_c
=\boxed{\frac{5\sh}{4m_er_c}}.
\end{equation}
With $r_c\simeq2.818\times10^{-15}$ m,
\begin{equation}
\boxed{v_{\rm equator}\simeq5.1\times10^{10}
\ {\rm m/s}\simeq171\,c}.
\end{equation}
This impossible speed shows that electron
spin is intrinsic angular momentum, not the
classical rotation of a tiny extended sphere.

\problem{4.29}{Spin commutators and Pauli-product identity}
For spin one-half, $S_j=\sh\sigma_j/2$,
where
\begin{equation}
\sigma_x=\begin{pmatrix}0&1\\1&0\end{pmatrix},
\quad\sigma_y=\begin{pmatrix}0&-\ii\\\ii&0\end{pmatrix},
\quad\sigma_z=\begin{pmatrix}1&0\\0&-1\end{pmatrix}.
\end{equation}
Explicit multiplication gives
$\sigma_x\sigma_y=\ii\sigma_z$,
$\sigma_y\sigma_x=-\ii\sigma_z$,
and $\sigma_j^2=I$. Cyclic permutation yields
\begin{equation}
\boxed{\sigma_j\sigma_k=
\delta_{jk}I+\ii\sum_\ell
\epsilon_{jk\ell}\sigma_\ell},
\end{equation}
hence
$\boxed{[S_j,S_k]=\ii\sh\epsilon_{jk\ell}S_\ell}$.

\problem{4.30}{Measured spins and uncertainties for $(3i,4)$}
The state $\chi=A(3\ii,4)^{\mathsf T}$
has $A=1/5$ (real, positive).
With the Pauli matrices,
\begin{equation}
\boxed{\avg{S_x}=0,\quad
\avg{S_y}=-\frac{12\sh}{25},\quad
\avg{S_z}=-\frac{7\sh}{50}}.
\end{equation}
Since $S_j^2=\sh^2I/4$,
\begin{equation}
\boxed{\sigma_{S_x}=\frac\sh2,\quad
\sigma_{S_y}=\frac{7\sh}{50},\quad
\sigma_{S_z}=\frac{12\sh}{25}}.
\end{equation}
The products
$\sigma_{S_x}\sigma_{S_y}=7\sh^2/100$ and
$\sigma_{S_x}\sigma_{S_z}=6\sh^2/25$
equal $(\sh/2)|\avg{S_z}|$
and $(\sh/2)|\avg{S_y}|$, respectively.
The remaining inequality is automatic because
$\avg{S_x}=0$.

\problem{4.31}{Expectations of a general spin-one-half state}
For normalized $\chi=(a,b)^{\mathsf T}$,
$|a|^2+|b|^2=1$, matrix multiplication yields
\begin{equation}
\boxed{\avg{S_x}=\sh\,\mathrm{Re}(a^*b),\quad
\avg{S_y}=\sh\,\mathrm{Im}(a^*b),\quad
\avg{S_z}=\frac\sh2(|a|^2-|b|^2)}.
\end{equation}
Since all Pauli matrices square to $I$,
\begin{equation}
\boxed{\avg{S_x^2}=\avg{S_y^2}=\avg{S_z^2}
=\frac{\sh^2}{4}},\qquad
\boxed{\avg{S^2}=\frac{3\sh^2}{4}}.
\end{equation}
The spin standard deviations follow by
subtracting squared means from $\sh^2/4$.

\problem{4.32}{Eigenstates and probabilities for $S_y$}
Diagonalize
$S_y=(\sh/2)\begin{pmatrix}0&-\ii\\\ii&0\end{pmatrix}$:
\begin{equation}
\boxed{\lambda_\pm=\pm\sh/2,\qquad
\chi_{y,+}=\frac1{\sqrt2}\binom{1}{\ii},
\quad\chi_{y,-}=\frac1{\sqrt2}\binom{1}{-\ii}}.
\end{equation}
For a general normalized $\chi=(a,b)^T$,
\begin{equation}
\boxed{P_y(+)=\frac12|a-\ii b|^2,\quad
P_y(-)=\frac12|a+\ii b|^2}.
\end{equation}
Their sum is $|a|^2+|b|^2=1$.
Since $S_y^2=(\sh^2/4)I$, measuring
$S_y^2$ yields only $\boxed{\sh^2/4}$,
with probability $1$.

\problem{4.33}{Spin along an arbitrary spatial direction}
For unit $\hat{\mathbf n}=
(\sin\theta\cos\phi,\sin\theta\sin\phi,\cos\theta)$,
\begin{equation}
S_{\mathbf n}=\frac\sh2
\begin{pmatrix}
\cos\theta&\ee^{-\ii\phi}\sin\theta\\
\ee^{\ii\phi}\sin\theta&-\cos\theta
\end{pmatrix}.
\end{equation}
Because $(\boldsymbol\sigma\cdot\hat{\mathbf n})^2=I$,
the eigenvalues are $\pm\sh/2$.
Solving the first row using
$\cos\theta\mp1=\mp2\sin^2(\theta/2)$
gives normalized eigenspinors
\begin{equation}
\boxed{\chi_{\mathbf n,+}=
\binom{\cos(\theta/2)}
{\ee^{\ii\phi}\sin(\theta/2)}},
\quad
\boxed{\chi_{\mathbf n,-}=
\binom{\ee^{-\ii\phi}\sin(\theta/2)}
{-\cos(\theta/2)}}.
\end{equation}
They are orthogonal and have unit norm.

\problem{4.34}{Spin-one matrices}
Choose the ordered $S_z$ basis
$\ket{1,1}$, $\ket{1,0}$, $\ket{1,-1}$.
The ladder rule
$S_\pm\ket{1,m}=\sh
\sqrt{2-m(m\pm1)}\ket{1,m\pm1}$
gives
\begin{align}
\boxed{S_x}&=\frac\sh{\sqrt2}
\begin{pmatrix}0&1&0\\1&0&1\\0&1&0\end{pmatrix},\\
\boxed{S_y}&=\frac\sh{\sqrt2}
\begin{pmatrix}0&-\ii&0\\\ii&0&-\ii\\0&\ii&0\end{pmatrix},\\
\boxed{S_z}&=\sh\begin{pmatrix}
1&0&0\\0&0&0\\0&0&-1
\end{pmatrix}.
\end{align}
Direct products give
$[S_x,S_y]=\ii\sh S_z$,
and cyclic permutations, and
$S_x^2+S_y^2+S_z^2=2\sh^2I_3$.
